% Chapter 15: Hierarchical linear models
% Bayesian Data Analysis - Gelman et al.

\chapter{分层线性模型}
\label{ch:hlm}

\section{本章导论}
\label{sec:ch15-intro}

在第七章和第十二章中，我们已经掌握了贝叶斯推断的基本框架，并学会了使用马尔科夫链蒙特卡洛（MCMC）方法处理复杂的后验分布。当数据具有分组结构时——例如，来自多所学校的学生成绩、多次临床试验的患者反应、或不同工厂的生产质量指标——我们需要一种更精细的工具来同时捕捉组内相关性和组间异质性。

分层线性模型（Hierarchical Linear Models, HLM）正是为解决这个问题而诞生的统计框架。它允许我们在不同层次上对数据进行建模：个体层次的变异和组层次的变异自然地被纳入统一的概率模型中。从历史上看，分层线性模型的发展源于教育研究（中学术的"学校效应"问题）、生物统计（生长曲线分析）和社会学（邻里效应研究），并在 1980 年代由 Bryk, Raudenbush 等人系统化为一套完整的统计方法论。

本章我们将学习如何将分层思想与线性回归结合，建立分层线性模型；如何用贝叶斯方法进行推断；以及如何从矩阵形式理解模型的本质。我们将看到，分层模型不仅是处理分组数据的工具，更是一种在"合并信息"与"保持特异性"之间寻求平衡的哲学。

\section{从分组数据到分层模型}
\label{sec:grouped-to-hierarchical}

\subsection{传统方法的困境}
\label{sec:traditional-dilemma}

考虑一个经典问题：估计 $J$ 所学校的培训效果。每所学校 $j$ 报告了一个效果估计 $y_j$ 及其标准误 $\sigma_j$。一个 naive 的做法是简单地将所有 $y_j$ 取平均作为总体效果估计。但这种方法忽略了一个关键事实：不同学校之间本身就可能存在系统性差异。

另一方面，如果我们分别为每所学校估计效果，就完全忽略了数据可以"相互学习"这一事实。例如，学校 3 的估计精度很低（$\sigma_3$ 很大），但其他学校的数据可以帮助我们改进对它的估计。

这就是分层模型要解决的核心问题：如何在保持组特异性估计的同时，充分利用数据的层次结构进行信息共享。

\subsection{两层次模型框架}
\label{sec:two-level-framework}

一般的两层次线性模型可以写为：

\begin{enumerate}
  \item \textbf{第一层（组内模型）}：$y_{ij} = \mathbf{x}_{ij}^\top \boldsymbol{\beta}_j + \epsilon_{ij}$，其中 $\epsilon_{ij} \sim \mathrm{N}(0, \sigma^2)$
  \item \textbf{第二层（组间模型）}：$\boldsymbol{\beta}_j = \mathbf{A}_j \boldsymbol{\gamma} + \mathbf{u}_j$，其中 $\mathbf{u}_j \sim \mathrm{N}(\mathbf{0}, \boldsymbol{\Sigma})$
\end{enumerate}

这里 $y_{ij}$ 是第 $j$ 组第 $i$ 个观测值，$\mathbf{x}_{ij}$ 是协变量向量，$\boldsymbol{\beta}_j$ 是第 $j$ 组的效应向量，$\boldsymbol{\gamma}$ 是总体（总体层面）效应，$\mathbf{A}_j$ 是已知的设计矩阵，$\boldsymbol{\Sigma}$ 是组间协方差矩阵。

将两层结合，我们可以写出完整的分层模型：

\[
y_{ij} = \mathbf{x}_{ij}^\top \mathbf{A}_j \boldsymbol{\gamma} + \mathbf{x}_{ij}^\top \mathbf{u}_j + \epsilon_{ij}
\]

其中 $\mathbf{u}_j$ 和 $\epsilon_{ij}$ 相互独立。

\section{分层线性模型的矩阵形式}
\label{sec:matrix-formulation}

\subsection{模型的一般表达}
\label{sec:general-matrix-form}

将所有数据堆叠为向量形式，分层线性模型可以紧凑地写为：

\[
\mathbf{y} = \mathbf{X}\boldsymbol{\gamma} + \mathbf{Z}\mathbf{u} + \boldsymbol{\epsilon}
\]

其中：
\begin{itemize}
  \item $\mathbf{y}$ 是 $n \times 1$ 观测向量
  \item $\mathbf{X}$ 是 $n \times p$ 固定效应设计矩阵
  \item $\boldsymbol{\gamma}$ 是 $p \times 1$ 固定效应参数
  \item $\mathbf{Z}$ 是 $n \times q$ 随机效应设计矩阵
  \item $\mathbf{u}$ 是 $q \times 1$ 随机效应向量，$\mathbf{u} \sim \mathrm{N}(\mathbf{0}, \boldsymbol{\Sigma})$
  \item $\boldsymbol{\epsilon}$ 是 $n \times 1$ 误差向量，$\boldsymbol{\epsilon} \sim \mathrm{N}(\mathbf{0}, \sigma^2 \mathbf{I}_n)$
\end{itemize}

随机效应 $\mathbf{u}$ 和误差 $\boldsymbol{\epsilon}$ 相互独立，因此 $\mathbf{y}$ 的边际分布为：

\[
\mathbf{y} \sim \mathrm{N}(\mathbf{X}\boldsymbol{\gamma}, \mathbf{V}), \quad \mathbf{V} = \sigma^2 \mathbf{I}_n + \mathbf{Z}\boldsymbol{\Sigma}\mathbf{Z}^\top
\]

\begin{Remark}
边际分布 $\mathbf{y} \sim \mathrm{N}(\mathbf{X}\boldsymbol{\gamma}, \mathbf{V})$ 正是我们在第八章学到的贝叶斯线性回归形式，只不过这里的协方差矩阵 $\mathbf{V}$ 依赖于方差参数 $\sigma^2$ 和 $\boldsymbol{\Sigma}$。
\end{Remark}

\subsection{组结构的矩阵表示}
\label{sec:group-structure}

考虑 $J$ 组、每组 $n_j$ 观测的经典情形。设 $\mathbf{y}_j$ 是第 $j$ 组的 $n_j \times 1$ 观测向量，$\mathbf{X}_j$ 是对应的设计矩阵。模型可写为：

\[
\mathbf{y}_j = \mathbf{X}_j \boldsymbol{\gamma} + \mathbf{Z}_j \mathbf{u}_j + \boldsymbol{\epsilon}_j, \quad \boldsymbol{\epsilon}_j \sim \mathrm{N}(\mathbf{0}, \sigma^2 \mathbf{I}_{n_j})
\]

其中 $\mathbf{u}_j \sim \mathrm{N}(\mathbf{0}, \boldsymbol{\Sigma}_\mathbf{u})$ 是组 $j$ 的随机效应。

整个数据的似然函数为：

\[
p(\mathbf{y} \mid \boldsymbol{\gamma}, \sigma^2, \boldsymbol{\Sigma}) = \prod_{j=1}^J \mathrm{N}(\mathbf{y}_j \mid \mathbf{X}_j\boldsymbol{\gamma}, \sigma^2 \mathbf{I}_{n_j} + \mathbf{Z}_j \boldsymbol{\Sigma} \mathbf{Z}_j^\top)
\]

\section{后验推断}
\label{sec:posterior-inference-hlm}

\subsection{完全条件分布}
\label{sec:full-conditional}

在贝叶斯分层模型中，我们通常关注以下后验分布：

\[
p(\boldsymbol{\gamma}, \mathbf{u}, \sigma^2, \boldsymbol{\Sigma} \mid \mathbf{y})
\]

为方便 MCMC 抽样，我们使用 Data augmentation 技术，引入隐变量 $\mathbf{u}$。完全数据模型为：

\[
\mathbf{y}_j \mid \mathbf{u}_j, \boldsymbol{\gamma}, \sigma^2 \sim \mathrm{N}(\mathbf{X}_j\boldsymbol{\gamma} + \mathbf{Z}_j\mathbf{u}_j, \sigma^2 \mathbf{I}_{n_j})
\]
\[
\mathbf{u}_j \mid \boldsymbol{\Sigma} \sim \mathrm{N}(\mathbf{0}, \boldsymbol{\Sigma})
\]

完全条件分布为：

\[
\boldsymbol{\gamma} \mid \mathbf{y}, \mathbf{u}, \sigma^2 \sim \mathrm{N}(\boldsymbol{\mu}_\gamma, \boldsymbol{\Sigma}_\gamma)
\]
\[
\mathbf{u}_j \mid \mathbf{y}, \boldsymbol{\gamma}, \sigma^2, \boldsymbol{\Sigma} \sim \mathrm{N}(\boldsymbol{\mu}_{u_j}, \boldsymbol{\Sigma}_{u_j})
\]
\[
\sigma^2 \mid \mathbf{y}, \boldsymbol{\gamma}, \mathbf{u}, \boldsymbol{\Sigma} \sim \text{Inverse-}\chi^2(\nu_n, \sigma_n^2)
\]

其中后验均值和方差的具体形式依赖于数据的分组结构。\footnote{完全条件分布的详细推导见附录 \cref{sec:appendix-hlm-conditional}。}

\subsection{分层先验的构建}
\label{sec:hierarchical-prior}

对超参数 $\boldsymbol{\Sigma}$ 的先验选择是分层模型的关键。几种常见选择：

\begin{enumerate}
  \item \textbf{逆 Wishart 先验}：$\boldsymbol{\Sigma} \sim \inverse-Wishart(\nu_0, \boldsymbol{\Sigma}_0)$
  \item \textbf{独立逆 Gamma 先验}：对方差分量独立建模
  \item \textbf{弱信息先验}：使用半参数或数据驱动的先验
\end{enumerate}

当方差分量接近零时，逆 Wishart 先验可能导致后验分布出现拖尾问题。因此，实际中常使用对数尺度上的先验，例如：

\[
\log(\sigma_j^2) \sim \mathrm{N}(\mu_0, \tau_0^2)
\]

\section{预测}
\label{sec:prediction-hlm}

分层模型的一个重要应用是预测新组或新观测的值。

\subsection{组均值预测}
\label{sec:group-mean-prediction}

给定新组 $j*$ 的协变量 $\mathbf{x}_{j*}$，我们希望预测其效应 $\boldsymbol{\beta}_{j*} = \mathbf{A}_{j*}\boldsymbol{\gamma} + \mathbf{u}_{j*}$。

后验预测分布为：

\[
\boldsymbol{\beta}_{j*} \mid \mathbf{y} \sim \mathrm{N}(\boldsymbol{\mu}_{\beta_{j*}}, \boldsymbol{\Sigma}_{\beta_{j*}})
\]

其中后验均值是总体效应 $\boldsymbol{\gamma}$ 的后验均值和组特异偏差 $\mathbf{u}_j$ 的后验均值的加权组合。\footnote{组均值预测的详细推导见附录 \cref{sec:appendix-hlm-prediction}。}

\subsection{新观测预测}
\label{sec:new-observation-prediction}

对于新组 $j*$ 的新观测 $y_{new}$，预测分布为：

\[
y_{new} \mid \mathbf{y} \sim \mathrm{N}(\mathbf{x}_{new}^\top \boldsymbol{\mu}_{\beta_{j*}}, \sigma^2 + \mathbf{x}_{new}^\top \boldsymbol{\Sigma}_{\beta_{j*}} \mathbf{x}_{new})
\]

注意预测方差包括两部分：观测误差方差 $\sigma^2$ 和参数不确定性方差 $\mathbf{x}_{new}^\top \boldsymbol{\Sigma}_{\beta_{j*}} \mathbf{x}_{new}$。

\section{方差分量估计}
\label{sec:variance-components}

\subsection{方差分量模型}
\label{sec:variance-component-model}

最简单的分层线性模型是所谓的\textbf{方差分量模型}（variance components model）：

\[
y_{ij} = \mu + u_j + \epsilon_{ij}
\]

其中 $u_j \sim \mathrm{N}(0, \sigma_u^2)$ 是组效应，$\epsilon_{ij} \sim \mathrm{N}(0, \sigma_\epsilon^2)$ 是个体误差。组内相关系数（ICC）为：

\[
\mathrm{ICC} = \frac{\sigma_u^2}{\sigma_u^2 + \sigma_\epsilon^2}
\]

ICC 衡量了总变异中有多大比例来自组间变异。

\subsection{边缘化与 EM 算法}
\label{sec:marginalization-em}

对随机效应 $\mathbf{u}$ 进行边缘化后，似然函数涉及高维积分。使用 EM 算法可以方便地估计方差分量。

E 步：计算随机效应的后验均值和协方差
\[
\widehat{\mathbf{u}}_j^{(t)} = \boldsymbol{\Sigma}_j \mathbf{Z}_j^\top \mathbf{V}_j^{-1} (\mathbf{y}_j - \mathbf{X}_j \widehat{\boldsymbol{\gamma}}^{(t)})
\]

M 步：更新参数估计
\[
\widehat{\boldsymbol{\gamma}}^{(t+1)} = \left(\sum_j \mathbf{X}_j^\top \mathbf{V}_j^{-1} \mathbf{X}_j \right)^{-1} \sum_j \mathbf{X}_j^\top \mathbf{V}_j^{-1} (\mathbf{y}_j - \mathbf{Z}_j \widehat{\mathbf{u}}_j^{(t)})
\]

方差分量的更新涉及更复杂的估计方程。\footnote{EM 算法对方差分量的更新公式见附录 \cref{sec:appendix-em-variance}。}

\section{相关作业干预示例}
\label{sec:example- coaching}

考虑一个经典的教育研究问题：估计 $J = 8$ 所学校的培训效果。

\subsection{问题背景}
\label{sec:coaching-background}

每所学校报告了培训效果估计 $y_j$ 和标准误 $\sigma_j$。数据如下：

\begin{center}
\begin{tabular}{ccc}
\hline
学校 $j$ & 估计 $y_j$ & 标准误 $\sigma_j$ \\
\hline
1 & 28.4 & 14.9 \\
2 & 7.0 & 10.2 \\
3 & -4.4 & 16.3 \\
4 & 7.0 & 11.0 \\
5 & 1.9 & 9.4 \\
6 & 6.6 & 10.5 \\
7 & 10.8 & 17.6 \\
8 & 33.1 & 15.6 \\
\hline
\end{tabular}
\end{center}

\subsection{模型设定}
\label{sec:coaching-model}

我们设定分层模型：

\[
y_j \mid \theta_j \sim \mathrm{N}(\theta_j, \sigma_j^2)
\]
\[
\theta_j \mid \mu, \tau \sim \mathrm{N}(\mu, \tau^2)
\]
\[
\mu \sim \mathrm{N}(0, 100^2), \quad \tau \sim \mathrm{Half-Cauchy}(0, 25)
\]

其中 $\theta_j$ 是学校 $j$ 的真实效果，$\mu$ 是总体效果，$\tau$ 是学校间标准差。

\subsection{后验推断}
\label{sec:coaching-posterior}

使用 MCMC 方法，我们可以得到 $\theta_j$ 的后验分布。关键发现是：

\begin{Remark}
后验均值是先验均值（0）和数据均值 $y_j$ 的加权平均，权重由精度（方差的倒数）给出。当 $\tau$ 小时，所有 $\theta_j$ 被强烈收缩到公共均值 $\mu$；当 $\tau$ 大时，每个 $\theta_j$ 主要由其自身数据决定。
\end{Remark}

这个例子展示了分层模型的核心优势：它提供了在\textbf{合并信息}和\textbf{保持特异性}之间的自然权衡。

\section{稳健推断}
\label{sec:robust-inference}

当数据存在异常值或偏离正态假设时，标准分层线性模型可能失效。

\subsection{稳健误差分布}
\label{sec:robust-error}

一种方法是使用 Student-$t$ 分布替代正态分布作为误差分布：

\[
\epsilon_{ij} \sim \text{Student-}t(\nu, 0, \sigma^2)
\]

当 $\nu \to \infty$ 时，Student-$t$ 分布趋近于正态分布；当 $\nu$ 较小时，分布的尾部更重，对异常值更稳健。

另一种方法是使用污染正态模型：

\[
\epsilon_{ij} \sim (1 - \epsilon) \mathrm{N}(0, \sigma^2) + \epsilon \mathrm{N}(0, k\sigma^2)
\]

其中 $\epsilon$ 是污染比例，$k$ 是方差放大因子。

\subsection{分层先验的稳健性}
\label{sec:robust-prior}

对方差分量使用弱信息先验可以提高稳健性。Half-t 先验和 Horseshoe 先验是常用的选择。

\begin{Remark}
选择先验时，应该考虑参数的尺度。例如，在 $\log(\tau)$ 上指定先验通常比在 $\tau$ 上更合理，因为前者保证 $\tau$ 始终为正。
\end{Remark}

\section{本章小结}
\label{sec:ch15-summary}

本章我们学习了分层线性模型的贝叶斯分析框架：

\begin{enumerate}
  \item \textbf{模型构建}：通过两层结构同时捕捉组内变异和组间变异
  \item \textbf{矩阵形式}：将分层模型表达为 $\mathbf{y} = \mathbf{X}\boldsymbol{\gamma} + \mathbf{Z}\mathbf{u} + \boldsymbol{\epsilon}$
  \item \textbf{后验推断}：利用完全条件分布进行 MCMC 抽样
  \item \textbf{预测}：提供组均值和新观测的预测分布
  \item \textbf{方差分量}：通过 ICC 等指标量化组间变异的重要性
  \item \textbf{稳健推断}：使用 Student-$t$ 误差或污染正态模型处理异常值
\end{enumerate}

分层线性模型是现代贝叶斯统计中最重要的工具之一，它的核心思想——在部分合并和完全分离之间寻求平衡——贯穿于整个统计学。

\section{下节预告}
\label{sec:ch15-next}

下一章我们将学习广义线性模型（Generalized Linear Models），这是将分层框架推广到非正态数据的关键步骤。

\endinput
